\Needspace{9\baselineskip}\begin{center}\textit{Lösung der Kongruenz \(xx\equiv A\) durch das Ausschlussverfahren.}\end{center}


\Needspace{6\baselineskip}
\Article{319}
Die \textit{Lösbarkeit} der Kongruenz \(xx\equiv A\; (\mathrm{mod}.\ m)\), die mit der unbestimmten Gleichung \(xx=A+my\) übereinstimmt, haben wir in Abschnitt IV (Artikel \(146\)) so behandelt, dass nichts weiter zu wünschen übrig zu bleiben scheint. Was jedoch die Ermittlung der Unbekannten selbst betrifft, haben wir schon oben (Artikel \(152\)) bemerkt, dass indirekte Verfahren den direkten bei weitem vorzuziehen sind. Ist \(m\) eine Primzahl, auf welchen Fall sich die übrigen leicht zurückführen lassen, so könnten wir zu diesem Zweck die Indextafel I, nach der Bemerkung in Artikel \(316\) mit III verbunden, verwenden, wie wir in Artikel \(60\) allgemeiner gezeigt haben; doch wäre dieses Verfahren auf den Umfang der Tafel beschränkt. Aus diesen Gründen hoffen wir, dass das folgende allgemeine und bequeme Verfahren den Liebhabern der Arithmetik willkommen sein wird.

Zunächst bemerken wir, dass es genügt, nur diejenigen Werte von \(x\) zu ermitteln, die positiv und nicht größer als \(\frac12m\) sind, da jeder andere Wert nach dem Modul \(m\) einem von ihnen oder dessen Negativem kongruent ist. Bei einem solchen Wert von \(x\) muss der Wert von \(y\) zwischen den Grenzen \(-\frac Am\) und \(\frac14m-\frac Am\) liegen. Das Verfahren, das sich unmittelbar anbietet, bestünde also darin, für jeden innerhalb dieser Grenzen liegenden Wert von \(y\) (die Gesamtheit dieser Werte wollen wir mit \(\Omega\) bezeichnen) den Wert von \(A+my\), den wir mit \(V\) bezeichnen, zu berechnen und nur diejenigen beizubehalten, für die \(V\) ein Quadrat wird. Ist \(m\) eine kleine Zahl (\textit{z. B.} unter \(40\)), so ist dieser Versuch so kurz, dass er kaum einer Abkürzung bedarf; ist dagegen \(m\) groß, so lässt sich die Arbeit durch das folgende \textit{Ausschlussverfahren} beliebig verkürzen.

\Needspace{6\baselineskip}
\Article{320}
Es sei \(E\) eine beliebige ganze, zu \(m\) teilerfremde Zahl größer als \(2\), und alle ihre verschiedenen quadratischen Nichtreste (\textit{d. h.} nach \(E\) inkongruenten Nichtreste) seien \(a,\allowbreak b,\allowbreak c\) usw.;
\SourceLeaf{394}
schließlich seien die Wurzeln der Kongruenzen
\[A+my\equiv a,\quad A+my\equiv b,\quad A+my\equiv c\quad\text{usw. nach dem Modul }E\]
die Zahlen \(\alpha,\allowbreak \beta,\allowbreak \gamma\) usw., die alle positiv und kleiner als \(E\) gewählt werden dürfen. Gibt man also \(y\) einen Wert, der einer der Zahlen \(\alpha,\allowbreak \beta,\allowbreak \gamma\) usw. nach \(E\) kongruent ist, so ist der daraus entstehende Wert von \(V=A+my\) einer der Zahlen \(a,\allowbreak b,\allowbreak c\) usw. kongruent, also ein Nichtrest von \(E\), und kann daher kein Quadrat sein. Somit können aus \(\Omega\) sofort alle Zahlen der Formen \(Et+\alpha,\allowbreak Et+\beta,\allowbreak Et+\gamma\) usw. als unbrauchbar ausgeschlossen werden, und es genügt, die übrigen, deren Gesamtheit \(\Omega'\) sei, zu erproben. Bei dieser Operation kann man der Zahl \(E\) den Namen \textit{Ausschlusszahl} geben.

Nimmt man nun eine andere geeignete Zahl \(E'\) als Ausschlusszahl, so findet man auf genau dieselbe Weise ebenso viele Zahlen \(\alpha',\allowbreak \beta',\allowbreak \gamma'\) usw., wie sie verschiedene quadratische Nichtreste hat, denen \(y\) nach dem Modul \(E'\) nicht kongruent sein kann. Daher darf man wiederum aus \(\Omega'\) alle Zahlen der Formen \(E't+\alpha',\allowbreak E't+\beta',\allowbreak E't+\gamma'\) usw. entfernen. So kann man mit immer weiteren Ausschlusszahlen fortfahren, bis die Menge der Zahlen aus \(\Omega\) so weit verkleinert ist, dass es nicht schwieriger erscheint, alle übrigen tatsächlich zu erproben, als neue Ausschlüsse vorzunehmen.

\textit{Beispiel.} Ist die Gleichung \(xx=22+97y\) gegeben, so sind die Grenzen der Werte von \(y\) gleich \(-\frac{22}{97}\) und \(24\frac14-\frac{22}{97}\). Da der Wert \(0\) offenbar unbrauchbar ist, umfasst \(\Omega\) also die Zahlen \(1,\allowbreak 2,\allowbreak 3\ldots24\). Für \(E=3\) gibt es nur den einen Nichtrest \(a=2\); daraus folgt \(\alpha=1\). Somit sind aus \(\Omega\) alle Zahlen der Form \(3t+1\) auszuschließen; die Anzahl der verbleibenden Zahlen \(\Omega'\) ist \(16\). Ebenso hat man für \(E=4\) die Werte \(a=2,\allowbreak b=3\), also \(\alpha=0,\allowbreak \beta=1\); daher müssen die Zahlen der Formen \(4t\) und \(4t+1\) verworfen werden, und es bleiben diese acht: \(2,\allowbreak 3,\allowbreak 6,\allowbreak 11,\allowbreak 14,\allowbreak 15,\allowbreak 18,\allowbreak 23\). Ebenso findet man für \(E=5\), dass die Zahlen der Formen \(5t\) und \(5t+3\) zu verwerfen sind; es bleiben \(2,\allowbreak 6,\allowbreak 11,\allowbreak 14\). Die Ausschlusszahl \(6\) würde die Zahlen der Formen \(6t+1\) und \(6t+4\) entfernen; diese, die mit den Zahlen der Form \(3t+1\) übereinstimmen, fehlen aber bereits. Die Ausschlusszahl \(7\) entfernt die Zahlen der Formen \(7t+2,\allowbreak 7t+3,\allowbreak 7t+5\) und lässt \(6,\allowbreak 11,\allowbreak 14\) übrig. Setzt man diese für \(y\) ein, so erhält man der Reihe nach \(V=604,\allowbreak 1089,\allowbreak 1380\), wovon allein der zweite Wert ein Quadrat ist; daraus folgt \(x=\mp33\).
\SourceLeaf{395}

\Needspace{6\baselineskip}
\Article{321}
Die mit der Ausschlusszahl \(E\) ausgeführte Operation entfernt unter den Werten von \(V\), die den Werten von \(y\) in \(\Omega\) entsprechen, alle quadratischen Nichtreste von \(E\), lässt aber die Reste derselben Zahl unberührt. Daher ist leicht einzusehen, dass sich die Verwendung der Ausschlusszahlen \(E\) und \(2E\) nicht unterscheidet, wenn \(E\) ungerade ist, da in diesem Fall \(E\) und \(2E\) dieselben Reste und Nichtreste haben. Werden also die Zahlen \(3,\allowbreak 4,\allowbreak 5\) usw. nacheinander als Ausschlusszahlen verwendet, so sind die Zahlen \(6,\allowbreak 10,\allowbreak 14\) usw., die das Doppelte ungerader Zahlen sind, als überflüssig zu überspringen. Ferner ist klar, dass die beiden mit den Ausschlusszahlen \(E,\allowbreak E'\) ausgeführten Operationen alle Werte von \(V\) entfernen, die Nichtreste beider Zahlen \(E,\allowbreak E'\) oder einer von ihnen sind, und diejenigen übrig lassen, die Reste beider sind. Haben nun \(E\) und \(E'\) keinen gemeinsamen Teiler, so sind alle entfernten Zahlen Nichtreste und alle übrigen Reste des Produkts \(EE'\). Daher bewirkt die Verwendung der Ausschlusszahl \(EE'\) in diesem Fall offenbar genau dasselbe wie die Verwendung der beiden Zahlen \(E,\allowbreak E'\) und ist somit nach diesen überflüssig. Man darf also auch alle Ausschlusszahlen überspringen, die sich in zwei teilerfremde Faktoren zerlegen lassen, und braucht nur Primzahlen, die \(m\) nicht teilen, oder Primzahlpotenzen zu verwenden. Schließlich ist klar, dass nach Verwendung der Ausschlusszahl \(p^\mu\), einer Potenz der Primzahl \(p\), die Ausschlusszahl \(p\) oder \(p^\nu\) bei \(\nu<\mu\) überflüssig wird. Da nämlich \(p^\mu\) unter den Werten von \(V\) nur seine Reste übrig gelassen hat, sind erst recht keine Nichtreste von \(p\) oder einer beliebigen niedrigeren Potenz \(p^\nu\) mehr vorhanden. Ist aber \(p\) oder \(p^\nu\) bereits vor \(p^\mu\) verwendet worden, so kann die letztere offenbar nur solche Werte von \(V\) entfernen, die zugleich Reste von \(p\) (oder \(p^\nu\)) und Nichtreste von \(p^\mu\) sind. Daher genügt es, nur solche Nichtreste von \(p^\mu\) für \(a,\allowbreak b,\allowbreak c\) usw. zu nehmen.

\Needspace{6\baselineskip}
\Article{322}
Die Berechnung der Zahlen \(\alpha,\allowbreak \beta,\allowbreak \gamma\) usw., die einer gegebenen Ausschlusszahl \(E\) entsprechen, wird durch die folgenden Bemerkungen erheblich verkürzt. Es seien \(\mathfrak A,\allowbreak \mathfrak B,\allowbreak \mathfrak C\) usw. Wurzeln der Kongruenzen \(my\equiv a,\allowbreak my\equiv b,\allowbreak my\equiv c\) usw. \((\mathrm{mod.}\ E)\) und \(k\) eine Wurzel von \(my\equiv-A\); dann gilt offenbar \(\alpha\equiv\mathfrak A+k,\allowbreak \beta\equiv\mathfrak B+k,\allowbreak \gamma\equiv\mathfrak C+k\) usw. Müsste man nun \(\mathfrak A,\allowbreak \mathfrak B,\allowbreak \mathfrak C\) usw. tatsächlich durch Lösung dieser Kongruenzen ermitteln, so wäre dieser Weg zur Bestimmung von \(\alpha,\allowbreak \beta,\allowbreak \gamma\) usw. gewiss nicht kürzer als der oben gezeigte; doch ist dies keineswegs nötig. Ist nämlich erstens \(E\) eine Primzahl und \(m\) ein quadratischer Rest von \(E\), so folgt aus Artikel \(98\), dass \(\mathfrak A,\allowbreak \mathfrak B,\allowbreak \mathfrak C\) usw.,
\SourceLeaf{396}
die die Werte der Ausdrücke \(\frac am,\allowbreak \frac bm,\allowbreak \frac cm\) usw. \((\mathrm{mod.}\ E)\) sind, verschiedene Nichtreste von \(E\) sind und daher mit \(\alpha,\allowbreak \beta,\allowbreak \gamma\) usw. vollständig übereinstimmen, wenn man von ihrer Reihenfolge absieht, die hier keine Rolle spielt. Ist dagegen unter derselben Voraussetzung \(m\) ein Nichtrest von \(E\), so stimmen die Zahlen \(\mathfrak A,\allowbreak \mathfrak B,\allowbreak \mathfrak C\) usw. mit sämtlichen quadratischen Resten außer \(0\) überein. Ist \(E\) das Quadrat einer ungeraden Primzahl, \(=pp\), und \(p\) bereits als Ausschlusszahl verwendet worden, so genügt es nach dem vorigen Artikel, für \(a,\allowbreak b,\allowbreak c\) usw. diejenigen Nichtreste von \(pp\) zu nehmen, die Reste von \(p\) sind, \textit{d. h.} die Zahlen \(p,\allowbreak 2p,\allowbreak 3p\ldots pp-p\) (also sämtliche Zahlen unter \(pp\) außer \(0\), die durch \(p\) teilbar sind). Daraus ist leicht einzusehen, dass sich für \(\mathfrak A,\allowbreak \mathfrak B,\allowbreak \mathfrak C\) usw. genau dieselben Zahlen ergeben müssen, nur in anderer Reihenfolge. Setzt man ebenso nach Verwendung der Ausschlusszahlen \(p\) und \(pp\) nun \(E=p^3\), so genügt es, für \(a,\allowbreak b,\allowbreak c\) usw. die Produkte der einzelnen Nichtreste von \(p\) mit \(pp\) zu nehmen. Dann ergeben sich für \(\mathfrak A,\allowbreak \mathfrak B,\allowbreak \mathfrak C\) usw. entweder dieselben Zahlen oder die Produkte von \(pp\) mit den einzelnen Resten von \(p\) außer \(0\), je nachdem, ob \(m\) Rest oder Nichtrest von \(p\) ist. Nimmt man allgemein für \(E\) eine beliebige Primzahlpotenz, etwa \(p^\mu\), nachdem alle niedrigeren Potenzen bereits verwendet worden sind, so ergeben sich für \(\mathfrak A,\allowbreak \mathfrak B,\allowbreak \mathfrak C\) usw. die Produkte von \(p^{\mu-1}\) entweder mit sämtlichen Zahlen kleiner als \(p\), stets außer \(0\), wenn \(\mu\) gerade ist; oder mit sämtlichen Nichtresten von \(p\) kleiner als \(p\), wenn \(\mu\) ungerade ist und \(mRp\) gilt; oder mit sämtlichen Resten, wenn \(mNp\) gilt. — Ist \(E=4\), also \(a=2,\allowbreak b=3\), so hat man für \(\mathfrak A,\allowbreak \mathfrak B\) entweder \(2\) und \(3\) oder \(2\) und \(1\), je nachdem, ob \(m\equiv1\) oder \(\equiv3\; (\mathrm{mod}.\ 4)\) gilt. Setzt man nach Verwendung der Ausschlusszahl \(4\) nun \(E=8\), so hat man \(\alpha=5\), woraus sich für \(\mathfrak A\) die Werte \(5,\allowbreak 7,\allowbreak 1,\allowbreak 3\) ergeben, je nachdem, ob \(m\equiv1,\allowbreak 3,\allowbreak 5,\allowbreak 7\; (\mathrm{mod}.\ 8)\) gilt. Ist nun allgemein \(E\) eine beliebige höhere Potenz von zwei, etwa \(2^\mu\), nachdem die niedrigeren bereits verwendet worden sind, so muss man \(a=2^{\mu-1},\allowbreak b=3.2^{\mu-2}\) setzen, wenn \(\mu\) gerade ist. Daraus folgt \(\mathfrak A=2^{\mu-1},\allowbreak \mathfrak B=3.2^{\mu-2}\) oder \(=2^{\mu-2}\), je nachdem, ob \(m\equiv1\) oder \(\equiv3\) gilt. Ist dagegen \(\mu\) ungerade, so muss man \(a=5.2^{\mu-3}\) setzen; dann ist \(\mathfrak A\) gleich dem Produkt der Zahl \(2^{\mu-3}\) mit \(5,\allowbreak 7,\allowbreak 1\) oder \(3\), je nachdem, ob \(m\equiv1,\allowbreak 3,\allowbreak 5\) oder \(7\; (\mathrm{mod}.\ 8)\) gilt.

Sachkundige werden sich im Übrigen leicht einen Apparat ausdenken, mit dem die unbrauchbaren Werte von \(y\) \textit{mechanisch} aus \(\Omega\) entfernt werden können, nachdem die Zahlen \(\alpha,\allowbreak \beta,\allowbreak \gamma\) usw. für so viele Ausschlusszahlen berechnet worden sind, wie nötig erscheinen; doch können wir hier weder darauf noch auf andere Hilfsmittel zur Verkürzung der Arbeit eingehen.
